\section{Experiments}

The experiment workflow has three steps: map observed foods to coordinates $(i,j,k)$ and set $T_{i,j,k}=1$ accordingly; measure occupancy $|\mathcal{S}|$ and collisions (multiple foods at the same $(i,j,k)$); and generate candidate foods for cells in $\mathcal{E}$.
Initial analysis reports which regions of the index set $\{1,\ldots,I\} \times \{1,\ldots,J\} \times \{1,\ldots,K\}$ are dense ($T_{i,j,k}=1$ for many $(i,j,k)$ in a neighborhood), sparse, or structurally inconsistent under the chosen ontology.
From a representation-learning perspective, this is a small-scale study of \term{combinatorial generalization} under hard axis constraints on $i$, $j$, and $k$.
Empirically, the workflow behaves like a controlled stress test of whether culinary naming conventions can survive contact with coordinate geometry.
Axis categories in Figure~\ref{fig:tensor-projections-speculated} and Figure~\ref{fig:tensor-3d-speculated} are ordered by occupancy-profile hierarchical clustering as described in the Methods section (applied to the indices $i$, $j$, $k$ along each axis).
\input{sections/generated/figure_omissions}

\paragraph{Review outcomes.}
All LLM-generated candidates exported for review were accepted and included in the real-only tensor views; the distribution of model confidence for those accepted candidates is shown in Figure~\ref{fig:confidence-histogram}.

We also run a targeted mini-experiment on foods whose names explicitly contain ``salad'' (for example, chicken salad, egg salad, and Snickers salad\footnote{Snickers salad is a dessert-style potluck dish, commonly made with chopped Snickers bars, tart apples, and whipped topping (often with pudding), and is widely associated with Upper Midwest potluck and church-supper traditions (especially Minnesota and Iowa), though no single point of origin is well documented.}), asking the model to classify each item by structural starch placement and thus to assign a cube index $i$ (and implicitly $j$, $k$).
Table~\ref{tab:salad-cube-rule} reports only the subset that violates the Cube Rule salad definition (no structural starch faces; i.e., the assigned $i$ does not match the ``salad'' morphology index).
As a secondary aside experiment, we test the hypothesis that most violations arise from discrete starch inclusions (for example, croutons, pasta pieces, or candy chunks), which shift morphology from the salad index $i_{\text{salad}}$ toward the Cube Rule ``nachos'' index $i_{\text{nachos}}$ (outer dish plus inner starch).
\IfFileExists{sections/generated/salad_cube_rule_table.tex}{%
  \input{sections/generated/salad_cube_rule_table}%
}{%
  \paragraph{Mini-experiment data pending.}
  Run \texttt{uv run python scripts/salad\_cube\_rule\_experiment.py} to generate \texttt{paper/sections/generated/salad\_cube\_rule\_table.tex}.
}

\begin{figure}[t]
\centering
\includegraphics[width=0.49\textwidth]{../results/figures/canonical_plus_llm_identified_plus_llm_speculated/tensor_cube_vs_protein.png}\hfill
\includegraphics[width=0.49\textwidth]{../results/figures/canonical_plus_llm_identified_plus_llm_speculated/tensor_cube_vs_starch.png}
\vspace{0.5em}
\includegraphics[width=0.49\textwidth]{../results/figures/canonical_plus_llm_identified_plus_llm_speculated/tensor_protein_vs_starch.png}
\caption{Pairwise projections of $T$: slices over $(i,j)$, $(i,k)$, and $(j,k)$. Black squares mark canonical $(i,j,k)$ with $T_{i,j,k}=1$, gray squares mark accepted real LLM-identified cells, and white squares with black outlines mark LLM speculative cells; cells in $\mathcal{E}$ are shown without glyphs.\omitSpeculatedSummary}
\label{fig:tensor-projections-speculated}
\end{figure}

\begin{figure}[t]
\centering
\includegraphics[width=0.76\textwidth]{../results/figures/canonical_plus_llm_identified_plus_llm_speculated/tensor_3d.png}
\caption{Static 2D rendering of the 3D occupancy tensor $T$. Black markers indicate canonical $(i,j,k) \in \mathcal{S}$, gray markers indicate accepted real LLM-identified cells, and white markers with black outlines indicate LLM speculative cells; $(i,j,k) \in \mathcal{E}$ are shown without glyphs.}
\label{fig:tensor-3d-speculated}
\end{figure}

\begin{figure}[t]
\centering
\includegraphics[width=0.35\textwidth]{../results/figures/confidence_histogram.png}
\input{sections/generated/confidence_histogram_caption}
\end{figure}
